\documentclass[../report.tex]{subfiles}
\begin{document}

Phụ lục này tóm tắt \file{README.md}. Yêu cầu phần mềm: Logisim-evolution 4.1.0 (cần Java 21) và Python 3.8 trở lên.
Assembler và bộ test không cần thư viện Python ngoài; riêng script sinh biểu đồ của báo cáo cần \texttt{matplotlib}.

\section{Mở project và chạy mô phỏng}

\begin{enumerate}
  \item Mở Logisim-evolution 4.1.0, chọn File $\to$ Open rồi chọn một trong hai file:\\
        \file{circ/RV32I\_SingleCycle.circ}\\
        \file{circ/RV32I\_Pipeline.circ}
  \item Ctrl+K cho clock tự chạy, chỉnh tốc độ ở Simulate $\to$ Tick Frequency. Ctrl+T chạy từng nửa chu kỳ,
        Ctrl+R reset.
  \item Chương trình \file{demo.s} đã nạp sẵn trong ROM. Khi chạy xong, TTY in \texttt{Hello RISC-V!}, LED 7 đoạn hiện
        \texttt{00E8}, 8 LED sáng và chân \sig{halt} lên 1.
  \item Quan sát: 32 thanh ghi hiện trên khối \sig{RegisterFile}; RAM xem trên khối \sig{DMEM}; ở bản pipeline,
        Pipeline Monitor hiện PC và lệnh ở từng tầng cùng \sig{Stall}, \sig{LoadUse}, \sig{E\_Redirect}, \sig{FwdA},
        \sig{FwdB}.
\end{enumerate}

\section{Nạp chương trình mới}

\begin{lstlisting}
python tools/rvasm.py programs/t08_sort.s -o build/t08_sort
\end{lstlisting}

Lệnh trên sinh \file{build/t08\_sort.text.hex}, \file{.data.hex} và \file{.lst}. Trong Logisim, chuột phải ROM
\sig{IMEM} $\to$ Load Image và chọn file \file{.text.hex}; nếu chương trình có \texttt{.data} thì làm tương tự với RAM
\sig{DMEM} và file \file{.data.hex}. Có thể ghi thẳng vào file mạch:

\begin{lstlisting}
python tools/rvasm.py programs/t08_sort.s -o build/t08_sort \
    --circ circ/RV32I_Pipeline.circ --circ-out build/sort_pipeline.circ
\end{lstlisting}

Với bản NH, thêm \texttt{-{}-nohazard} để assembler chèn \texttt{nop}.

\section{Chạy test}

\begin{lstlisting}
export LOGISIM_JAR=/path/to/logisim-evolution-4.1.0-all.jar
python tools/run_tests.py circ/RV32I_SingleCycle.circ
python tools/run_tests.py circ/RV32I_Pipeline.circ
python tools/run_tests.py circ/RV32I_Pipeline.circ --programs programs/fuzz
python tools/run_tests.py circ/extra/RV32I_Pipeline_NoHazard.circ --nohazard
java -jar $LOGISIM_JAR -w ALU tests/vectors/ALU.txt circ/RV32I_Pipeline.circ
python tools/compare.py programs/t06_hazard.s
\end{lstlisting}

Lệnh \texttt{-w} chạy test vector của một module và in \texttt{Passed: n, Failed: m}. Lệnh này không chạy được khi
thư mục chứa file mạch có file autosave của một phiên Logisim đang mở (\file{.RV32I\_Pipeline.circ.autosave}); khi
đó Logisim cần hỏi có khôi phục bản autosave không. Đóng Logisim hoặc chạy trên một bản sao của file mạch.

\section{Dựng lại báo cáo}

\begin{lstlisting}
python latex/scripts/gen_results.py
cd latex && latexmk -pdf report.tex
\end{lstlisting}

Script thứ nhất sinh lại các bảng số chu kỳ, giản đồ pipeline và biểu đồ trong \file{latex/generated/} từ golden
model. Lệnh thứ hai biên dịch báo cáo (pdfLaTeX và biber).

\section{Kịch bản demo}

\begin{enumerate}
  \item Mở \file{RV32I\_SingleCycle.circ}, giới thiệu sơ đồ khối. Chạy \file{demo.s}, chỉ vào \sig{RegisterFile}, TTY
        và LED 7 đoạn. Nạp \file{t11\_io.s}, click vào khối Keyboard, gõ chữ thường và xem TTY in lại chữ hoa.
  \item Mở \file{RV32I\_Pipeline.circ}, chạy từng chu kỳ tới một lần load-use: Pipeline Monitor cho \sig{Stall} = 1 và
        bubble \texttt{00000000} ở EX. Chạy tới một lần rẽ nhánh nhảy: \sig{E\_Redirect} = 1 và hai lệnh phía sau
        thành \texttt{00000000}.
  \item Chạy \file{run\_tests.py} trên terminal, cho thấy \texttt{ALL PASSED} với SC và PL.
  \item Chạy \file{run\_tests.py} trên bản NH không có \texttt{-{}-nohazard} để thấy FAIL, rồi chạy lại với
        \texttt{-{}-nohazard} để thấy PASS.
\end{enumerate}

\end{document}
